\begin{xltabular}{\textwidth}{| >{\bfseries}l | X |}

% --- INTESTAZIONE PRIMA PAGINA ---
\caption{Caso d'Uso: Aggiornamento dello Stato del Sistema} \label{tab:uc_update_state} \\
\hline
\rowcolor{gray!15} Campo & \textbf{Contenuto} \\ \hline
\endfirsthead

% --- INTESTAZIONE PAGINE SUCCESSIVE ---
\multicolumn{2}{c}{{\tablename\ \thetable{} -- Continua dalla pagina precedente}} \\
\hline
\rowcolor{gray!15} Campo & \textbf{Contenuto} \\ \hline
\endhead

% --- PIÈ DI PAGINA (se si spezza) ---
\hline
\multicolumn{2}{r}{\textit{Continua nella pagina successiva...}} \\
\endfoot

% --- PIÈ DI PAGINA FINALE ---
\endlastfoot

% --- CONTENUTO ---
Nome & Aggiornamento dello Stato del Sistema \\ \hline
Livello & Obiettivo di Sistema \\ \hline
Descrizione & Il \texttt{GridMonitor} è un thread \textit{daemon} avviato insieme all'applicazione, ed è l'unico componente che agisce di propria iniziativa: tutti gli altri servizi vengono invocati su richiesta. Ogni 5 secondi fa avanzare la simulazione. \\ \hline
Pre-condizioni & \parbox[t]{\linewidth}{%
L'applicazione è avviata e il \texttt{GridMonitor} è in esecuzione.\\
Il ciclo produce effetti osservabili se esiste almeno una sessione \texttt{ACTIVE} oppure una prenotazione \texttt{RESERVED} scaduta.
} \\ \hline
Post-condizioni & Le prenotazioni scadute sono tornate \texttt{ACTIVE}. Ogni sessione attiva ha batteria, kWh erogati e costo aggiornati, e l'importo del tick è stato scalato dal portafoglio del Conducente. Temperatura e carico di ogni trasformatore riflettono il calore generato nel tick. \\ \hline
Attori & Grid Monitor (Sistema) \\ \hline
Unit Test & \texttt{GridMonitorTest} \\ \hline
Flusso Principale &
\begin{enumerate}[leftmargin=*, nosep]
    \item Il timer del \texttt{GridMonitor} scatta, ogni 5 secondi.
    \item Il sistema libera le prenotazioni scadute (\texttt{<<include>>}): le colonnine \texttt{RESERVED} il cui \texttt{expiration\_timestamp} è ormai passato tornano \texttt{ACTIVE}.
    \item Per ogni sessione \texttt{ACTIVE} il sistema calcola l'energia erogata nel tick, pari alla potenza della colonnina moltiplicata per la frazione d'ora corrispondente a 5 secondi. \textbf{L'energia erogata non dipende dalla strategia}: la batteria avanza alla stessa velocità in entrambe le modalità, perché la potenza è una caratteristica fisica della colonnina.
    \item La strategia determina invece il \emph{costo} del tick e il \emph{calore} prodotto:
    \begin{itemize}[leftmargin=*, nosep]
        \item \textbf{Fast Charge:} il Conducente paga la tariffa piena, cioè il prezzo base della colonnina; la sessione genera il doppio del calore.
        \item \textbf{Eco Charge:} sullo stesso prezzo base è applicato un ulteriore sconto del 10\%; la sessione genera metà del calore.
    \end{itemize}
    \item Il sistema addebita il costo del tick al portafoglio del Conducente e salva sessione e saldo in un'unica transazione.
    \item Il sistema somma il calore prodotto da tutte le sessioni, lo raggruppa per trasformatore e lo applica a ciascuno di essi, che ne ricalcola temperatura e percentuale di carico.
    \item Il sistema esegue infine il controllo di qualità: se una colonnina ha una media inferiore a 2.0 con almeno 5 valutazioni, viene aperta una segnalazione (Tab.~\ref{tab:uc_rating_alert}).
\end{enumerate} \\ \hline
Flussi Alternativi &
\textbf{3. Batteria al 100\%:} \newline
3.1. Se il tick porta la batteria al 100\%, la sessione passa a \texttt{COMPLETED} e viene chiusa. \newline
3.2. Il sistema emette la ricevuta con il costo accumulato tick dopo tick e riporta la colonnina ad \texttt{ACTIVE}. \newline\newline
\textbf{5. Credito esaurito durante la ricarica:} \newline
5.1. Prima dell'addebito il sistema verifica che il saldo copra il costo del tick. \newline
5.2. Se non lo copre, il tick non viene addebitato: il portafoglio non può in nessun caso scendere sotto lo zero. \newline
5.3. La sessione viene chiusa d'autorità seguendo lo stesso percorso della conclusione normale: ricevuta emessa con il costo effettivamente accumulato e colonnina di nuovo \texttt{ACTIVE}. \newline
5.4. Il Conducente potrà avviare una nuova ricarica solo dopo aver ricaricato il portafoglio. \newline\newline
\textbf{6. Allarme Termico (\texttt{<<extend>>}):} \newline
6.1. Se applicando il calore la temperatura di un trasformatore supera i 90\,°C, il trasformatore notifica i propri \textit{observer} con l'evento \texttt{THERMAL\_ALERT}. \newline
6.2. Il \texttt{LoadBalancerService} porta \textbf{prima} tutte le colonnine di quel trasformatore allo stato \texttt{OVERLOADED}, in un'unica transazione atomica, impedendo l'avvio di nuove ricariche sul ramo. \newline
6.3. \textbf{Solo dopo} interrompe le sessioni ancora in corso: ciascuna passa a \texttt{INTERRUPTED\_THERMAL}, il Conducente viene rimborsato del costo dell'ultimo tick e allo Station Operator viene comunque accreditata la sua quota sull'energia realmente erogata. Ogni chiusura forzata ha una propria transazione indipendente: il fallimento di un rimborso non annulla quelli già andati a buon fine. \newline
6.4. Quando la temperatura ridiscende sotto i 70\,°C, il trasformatore emette l'evento \texttt{COOLING\_COMPLETE} e le colonnine rimaste \texttt{OVERLOADED} tornano \texttt{ACTIVE}. La distanza fra le due soglie evita che il sistema oscilli attorno a un unico valore. \\ \hline

\end{xltabular}
